% =============================================
% CONCLUSION ET PERSPECTIVES
% =============================================
\chapter*{Conclusion et Perspectives}
\addcontentsline{toc}{chapter}{Conclusion et Perspectives}
\thispagestyle{fancy}

Ce Projet de Fin d'Ann\'{e}e, inscrit dans le cadre de notre formation au sein de
l'\textbf{EMSI}, nous a offert l'opportunit\'{e} d'appr\'{e}hender des d\'{e}fis architecturaux et
techniques concrets. En d\'{e}veloppant l'application back-end \textbf{PT\_BACK}, notre
objectif principal \'{e}tait de concevoir une passerelle API robuste d\'{e}di\'{e}e \`{a} la gestion
compl\`{e}te et p\'{e}renne du cycle de vie des projets collaboratifs.

Du point de vue m\'{e}thodologique, l'approche \textbf{Agile} adopt\'{e}e a permis d'assurer un
rythme de d\'{e}veloppement soutenu. Apr\`{e}s avoir identifi\'{e} les acteurs centraux
(Manager, Chef de Projet et D\'{e}veloppeur) et \'{e}tabli l'ensemble des besoins,
nous avons \'{e}rig\'{e} la mod\'{e}lisation conceptuelle \`{a} retranscrire num\'{e}riquement dans le
framework.

Par la d\'{e}livrance finale des sp\'{e}cifications, la r\'{e}alisation accomplie du
backend met en v\'{e}rit\'{e} les forces du framework \textbf{Laravel 11} et de l'\'{e}cosyst\`{e}me PHP moderne :
\begin{itemize}[leftmargin=*, label=\textcolor{emsiRed}{\ding{118}}]
    \item L'orchestration r\'{e}gie des flux complexes : Assignations directes ou en masse,
    gestion de rôles.
    \item L'automatisation s\'{e}curis\'{e}e du "Scaffolding" avec son espace de pr\'{e}paration documentaire en dur pour chaque projet.
    \item L'authentification r\'{e}solument ext\'{e}rioris\'{e}e et standardis\'{e}e en \textbf{Keycloak} pour
    toute gouvernance Entreprise.
\end{itemize}

Ce PFA nous a en outre permis de perfectionner notre ma\^{i}trise de la programmation
orient\'{e}e objet et de consolider notre p\'{e}dagogie logicielle en ce qui concerne la r\'{e}daction d'un
code maintenable et s\'{e}curis\'{e}.

\vspace{0.8cm}
\noindent\textcolor{emsiDark}{\Large\bfseries Perspectives d'\'{e}volution}

Plusieurs enrichissements pourront imm\'{e}diatement parfaire l'\'{e}difice actuel :
\begin{enumerate}[leftmargin=*]
    \item \textbf{Clients Front-end et Mobile} : Le couplage de cette API \`{a} un v\'{e}ritable
    dashboard r\'{e}actif (React/Vue/Angular) et \`{a} une interface mobile pour les D\'{e}veloppeurs en d\'{e}placement.
    \item \textbf{Notifications et Real-Time} : Int\'{e}gration de \textit{WebSockets} ou de Firebase
    pour diffuser instantan\'{e}ment aux Chefs de Projet l'accomplissement d'une SLA ou la compl\'{e}tion d'une t\^{a}che sans forcer le rafraîchissement.
    \item \textbf{CI/CD Pipeline} : Mettre en place un processus de d\'{e}ploiement automatis\'{e}
    des conteneurs \textit{Docker} (sur serveurs cloud AWS, ou OVH).
\end{enumerate}

Ainsi ce modeste projet nous d\'{e}livre \`{a} son heure la satisfaction de la cl\'{e}ture
de cycle de notre dipl\^{o}me d'ing\'{e}nieur de l'\textbf{EMSI}.
